\chapter{Estado del arte}

En este capítulo se revisan los trabajos que han tratado de recuperar información lingüística a partir de registros cerebrales no invasivos. Esta revisión se hace con el objetivo de identificar la evolución de los métodos más próximos al problema abordado en este trabajo: la clasificación contextual de palabras, el aprendizaje de representaciones neuronales preentrenadas y la transferencia entre MEG y EEG.

La comparación entre trabajos debe realizarse con precaución. Como hemos explicado en la sección anterior, bajo la denominación de decodificación del lenguaje se agrupan tareas muy diferentes:
\begin{itemize}
\setlength{\itemsep}{0pt}
\setlength{\parsep}{0pt}
\setlength{\topsep}{2pt}
\setlength{\partopsep}{0pt}
\item distinguir unas pocas vocales imaginadas, 
\item identificar un segmento de audio entre varios candidatos, 
\item recuperar una palabra de un vocabulario cerrado o 
\item generar una frase completa. 
\end{itemize} 
También varían por otro lado:
\begin{itemize}
\setlength{\itemsep}{0pt}
\setlength{\parsep}{0pt}
\setlength{\topsep}{2pt}
\setlength{\partopsep}{0pt}
\item la modalidad de registro, 
\item el número de sujetos, 
\item la duración disponible por persona, 
\item la forma de presentar el estímulo y 
\item el criterio utilizado para separar entrenamiento y prueba. 
\end{itemize} Por ello, los resultados numéricos solo son comparables cuando se mantienen constantes estas condiciones.

\section{Decodificación no invasiva del lenguaje}

La investigación en decodificación no invasiva del lenguaje ha avanzado desde experimentos muy controlados hacia escenarios progresivamente más naturales. Los primeros trabajos se centraban principalmente en unidades aisladas y vocabularios pequeños, ya que este planteamiento facilita la repetición de estímulos y permite utilizar clasificadores convencionales. Los métodos recientes, en cambio, intentan trabajar con habla o lectura continua, relacionar la actividad cerebral con representaciones obtenidas mediante modelos de lenguaje y aprovechar el contexto de las palabras anteriores y posteriores.

En esta evolución pueden distinguirse tres niveles. El primero consiste en detectar información general, como la presencia de habla o la identidad de un fonema. El segundo busca clasificar palabras o segmentos dentro de un conjunto de candidatos. El tercero pretende reconstruir secuencias completas mediante un sistema \textit{brain-to-text}. Este último nivel es el más próximo a una interfaz de comunicación, aunque también es el más difícil de evaluar, porque un modelo de lenguaje externo puede generar frases plausibles incluso cuando la señal cerebral contiene poca información. En consecuencia, la clasificación de palabras continúa siendo una tarea intermedia útil para estudiar cuánto contenido lingüístico puede recuperarse realmente de la señal.

\subsection{Decodificación mediante EEG}

Una parte importante de la literatura basada en EEG se ha dedicado al habla imaginada. En estos experimentos se solicita a los participantes que imaginen la pronunciación de vocales, sílabas, palabras o comandos sin producir sonido. Las primeras propuestas combinaron características temporales, espectrales o tiempo-frecuencia con clasificadores como LDA, SVM o \textit{random forest}. Posteriormente se incorporaron redes convolucionales y recurrentes para aprender las características directamente a partir de la señal \cite{panachakel2021covert,alzahrani2024systematic}.

Los resultados obtenidos en este ámbito muestran que es posible diferenciar un número reducido de clases en condiciones controladas, especialmente cuando se entrena y evalúa un modelo específico para cada sujeto. Por ejemplo, se han utilizado CNNs para clasificar las cinco vocales imaginadas y distintos comandos direccionales, aunque el rendimiento cambia considerablemente entre participantes y conjuntos de datos \cite{tamm2020vowels,sarmiento2021vowels}. Estas diferencias dificultan extraer conclusiones generales. Un porcentaje elevado en una clasificación de dos o cinco clases no implica que el sistema pueda extenderse a un vocabulario amplio, ni que mantenga su funcionamiento con personas o sesiones no observadas durante el entrenamiento.

La lectura natural proporciona un escenario diferente. En este caso no se pide al participante que genere internamente una palabra concreta, sino que se registra su actividad mientras procesa un texto. El conjunto de datos ZuCo permitió avanzar en esta dirección al combinar EEG de alta densidad con seguimiento ocular. Las fijaciones indican qué palabra está observando el participante y permiten asociar intervalos de señal a elementos concretos de una frase \cite{hollenstein2020zuco2}. Esta alineación ha facilitado el desarrollo de modelos que reciben una secuencia de representaciones EEG a nivel de palabra e intentan reconstruir el texto leído.

Wang y Ji propusieron uno de los primeros sistemas de vocabulario abierto sobre ZuCo. Su método utiliza las representaciones EEG como entrada de un modelo basado en BART y plantea conjuntamente la traducción EEG-texto y la clasificación de sentimiento \cite{wang2022openvocab}. Más adelante, DeWave introdujo una codificación discreta de las ondas EEG antes de conectarlas con un decodificador lingüístico, con el objetivo de reducir la distancia entre una señal continua y la representación utilizada por un modelo de lenguaje \cite{duan2023dewave}. Estos trabajos fueron relevantes porque desplazaron el problema desde una clasificación cerrada hacia la generación de secuencias completas.

Sin embargo, una revisión posterior mostró que los resultados de varios sistemas EEG-to-text estaban afectados por la forma de realizar la evaluación. En particular, el uso implícito de \textit{teacher forcing} permitía que el decodificador recibiera palabras correctas de la secuencia durante la inferencia. Al eliminar esta ayuda y comparar el sistema con entradas de ruido, el rendimiento podía ser similar al obtenido sin información EEG \cite{jo2024eegtext}. Esta observación no invalida el uso de ZuCo ni la investigación sobre lectura natural, pero sí obliga a separar la capacidad lingüística del decodificador de la información que aporta realmente la señal cerebral.

Otro resultado importante procede del estudio de Sato et al., que reunió 175 horas de EEG de un único participante leyendo texto en voz alta. El modelo se entrenó mediante aprendizaje contrastivo entre segmentos EEG y audio y alcanzó un rendimiento muy superior al obtenido al limitar el entrenamiento a unas diez horas \cite{sato2024scaling}. El trabajo demuestra que la profundidad de datos por sujeto puede producir una mejora muy grande. No obstante, la tarea consiste en identificar segmentos de habla producida de cinco segundos y no palabras durante lectura silenciosa. Además, aunque se incorporan controles frente a la actividad muscular, el habla pronunciada introduce una fuente de información que no está presente en la percepción o en la lectura natural.

El trabajo de d'Ascoli et al. amplió la escala de la comparación al reunir registros EEG y MEG de cientos de participantes expuestos a millones de palabras en tareas de lectura y escucha \cite{dascoli2024words}. Sus resultados muestran que una arquitectura común puede recuperar representaciones de palabras en varios idiomas y dispositivos. Al mismo tiempo, el estudio señala dos diferencias relevantes para este trabajo: el rendimiento es mayor con MEG que con EEG y, manteniendo el contenido lingüístico, la lectura resulta más fácil de decodificar que la escucha. También se observa que la cantidad de datos disponible por participante tiene un efecto importante, lo que coincide con la tendencia encontrada en los registros EEG de larga duración.

En conjunto, la decodificación mediante EEG presenta dos líneas que todavía no se han unido de forma satisfactoria. Por una parte, existen clasificadores de habla imaginada capaces de distinguir vocabularios pequeños, pero suelen ser dependientes del sujeto y del protocolo. Por otra, los sistemas de lectura natural intentan generar texto abierto, aunque su evaluación puede estar dominada por el modelo lingüístico. Entre ambos extremos queda la clasificación contextual de palabras, que permite utilizar estímulos naturales y, al mismo tiempo, medir directamente la contribución de la señal cerebral. La disponibilidad de EEG, su menor coste y su portabilidad hacen que esta tarea resulte especialmente interesante, pero la baja relación señal-ruido y la heterogeneidad de los montajes continúan limitando su rendimiento.

\subsection{Decodificación mediante MEG}

La MEG se ha utilizado con menor frecuencia que EEG debido al coste y a la menor disponibilidad de los sistemas de adquisición. No obstante, la mayor calidad espacial de sus registros y la ausencia de una referencia eléctrica han favorecido resultados más consistentes en decodificación del lenguaje. Los primeros estudios, al igual que en EEG, se centraron en clasificaciones cerradas. Dash et al. analizaron cinco frases imaginadas y pronunciadas por ocho participantes. Mediante CNNs aplicadas a características espaciales, temporales y espectrales, obtuvieron resultados elevados dentro de ese conjunto reducido de frases \cite{dash2020megphrases}. El experimento mostró que la MEG contiene información discriminativa sobre la producción y la imaginación del habla, aunque el vocabulario cerrado y el número de participantes no permiten extrapolar esos porcentajes a una transcripción natural.

Un cambio de enfoque se produjo con la utilización de objetivos contrastivos. Défossez et al. entrenaron un codificador neuronal para acercar segmentos M/EEG a representaciones del audio escuchado \cite{defossez2023speech}. En lugar de asignar una etiqueta lingüística fija, el sistema debía identificar el fragmento de audio correspondiente entre varios candidatos. La arquitectura incorporaba atención espacial para trabajar con diferentes disposiciones de sensores, capas específicas de sujeto y convoluciones dilatadas. Este planteamiento demostró que las representaciones obtenidas por modelos de audio preentrenados constituyen objetivos útiles para la percepción del habla. Sin embargo, el método necesita disponer del audio candidato y no produce directamente palabras o texto.

La publicación de conjuntos de datos de habla continua ha permitido aumentar la escala de estos experimentos. Armeni et al. registraron aproximadamente diez horas por participante durante la escucha de una narración, mientras que MEG-MASC reunió historias naturales repetidas en varios sujetos \cite{armeni2022meg,gwilliams2023megmasc}. Más recientemente, LibriBrain ha proporcionado más de 52 horas de MEG de una sola persona escuchando audiolibros en inglés, con anotaciones temporales de habla, fonemas y palabras \cite{ozdogan2025libribrain}. Sus líneas base confirman que el rendimiento aumenta al incorporar más horas de entrenamiento y que los conjuntos profundos, aunque tengan pocos sujetos, son especialmente útiles para estudiar tareas lingüísticas complejas.

Sobre estos datos se ha consolidado la clasificación contextual de palabras. En lugar de procesar cada ventana de forma independiente, las representaciones neuronales asociadas a varias palabras consecutivas se introducen en un Transformer. El modelo puede así utilizar la estructura de la frase y las dependencias presentes en la señal cerebral. El trabajo de d'Ascoli et al. obtuvo una mejora clara respecto a un codificador local y estableció una referencia común basada en la recuperación de \textit{embeddings} lingüísticos \cite{dascoli2024words}. Esta propuesta se describe con mayor detalle más adelante, ya que constituye la tarea utilizada para ajustar y evaluar MEG-XL.

El paso desde palabras aisladas hacia secuencias completas se aborda en \textit{Unlocking Non-Invasive Brain-to-Text}. En este trabajo, las distribuciones de candidatos producidas por un clasificador de palabras se combinan con un modelo de lenguaje durante una búsqueda en haz. Además, se introduce un detector de palabras fuera del vocabulario y un mecanismo de rellenado contextual \cite{jayalath2025unlocking}. Los autores muestran que el reordenamiento lingüístico y la combinación selectiva de conjuntos de datos mejoran las métricas de secuencia y superan controles aleatorios más estrictos que los utilizados en algunos trabajos anteriores. Aun así, el sistema sigue dependiendo de un vocabulario de recuperación inicial y su calidad está limitada por la precisión del codificador neuronal.

La evolución de la decodificación mediante MEG muestra, por tanto, que el contexto y la cantidad de datos son dos factores fundamentales. Las ventanas locales permiten recuperar información acústica o léxica, pero las secuencias de palabras mejoran el resultado. Del mismo modo, los conjuntos con muchas horas por participante producen mejores modelos que aquellos que distribuyen el mismo tiempo entre un número muy elevado de personas. La principal desventaja es que la mayor parte de estos resultados se ha obtenido con equipos MEG concretos y en condiciones de laboratorio. Para que sus avances puedan trasladarse a una BCI más accesible es necesario determinar qué representaciones pueden conservarse al cambiar a EEG.

\section{Modelos preentrenados de señales cerebrales}

El aprendizaje supervisado utilizado en los primeros decodificadores obliga a disponer de una etiqueta para cada segmento. En lenguaje, esto requiere una sincronización precisa con palabras, fonemas o fragmentos de audio. El preentrenamiento autosupervisado intenta reducir esta dependencia aprendiendo primero regularidades generales de la señal y adaptando después el modelo a una tarea lingüística concreta. Esta estrategia ha dado lugar a los denominados modelos fundacionales de señales cerebrales.

La idea común consiste en reunir datos procedentes de diferentes sujetos y experimentos, dividir la señal en fragmentos o tokens y definir una tarea que no necesite anotaciones manuales. Entre los objetivos habituales se encuentran la reconstrucción de intervalos enmascarados, la predicción de representaciones latentes y la cuantización de la señal. El modelo preentrenado se ajusta después a tareas de clasificación, detección o regresión utilizando una cantidad menor de datos etiquetados. El éxito de esta transferencia depende de que las propiedades aprendidas durante el preentrenamiento sean útiles en el nuevo dominio.

\subsection{Modelos para EEG}

La mayoría de los modelos fundacionales iniciales se desarrolló para EEG, ya que existen más bases de datos públicas y la adquisición es más accesible. LaBraM utiliza predicción enmascarada de parches para aprender representaciones genéricas a partir de un corpus EEG de gran tamaño \cite{jiang2024labram}. EEGPT adopta igualmente un Transformer preentrenado, con el objetivo de conservar dependencias espaciales y temporales que puedan reutilizarse en distintos conjuntos y tareas \cite{wang2024eegpt}. Ambos trabajos muestran que el preentrenamiento puede mejorar determinados problemas de clasificación, aunque sus evaluaciones se concentran principalmente en aplicaciones clínicas y paradigmas BCI distintos del lenguaje continuo.

CBraMod introduce una arquitectura de atención \textit{criss-cross} para tratar por separado las relaciones temporales y las relaciones entre canales \cite{wang2025cbramod}. Esta factorización reduce el coste respecto a una atención completa y respeta la estructura bidimensional de la señal. El modelo resulta especialmente próximo a MEG-XL desde el punto de vista arquitectónico, pero fue diseñado y evaluado como un modelo general de EEG. Sus ventanas de preentrenamiento continúan siendo mucho más cortas que la escala temporal de una narración.

Otra línea se centra en aprender tokens discretos. BioCodec adapta los códecs neuronales a bioseñales y utiliza cuantización vectorial residual para comprimir cada canal en una secuencia de códigos \cite{avramidis2025biocodec}. El tokenizer puede reutilizarse como representación congelada o como objetivo de una tarea enmascarada. Esta separación entre tokenizer y modelo contextual permite que cada componente se entrene con datos distintos. De hecho, MEG-XL utiliza BioCodec, aunque el tokenizer original fue entrenado con EEG. Este resultado sugiere que algunas regularidades temporales de bajo nivel pueden compartirse entre modalidades, pero no demuestra por sí solo que un modelo contextual aprendido con MEG pueda transferirse a EEG.

BrainOmni aborda de forma explícita la heterogeneidad entre EEG y MEG. Su BrainTokenizer incorpora la posición, orientación y tipo de sensor para cuantizar registros procedentes de equipos distintos. A partir de estos tokens, el modelo se preentrena conjuntamente con miles de horas de EEG y cientos de horas de MEG \cite{xiao2025brainomni}. Esta propuesta constituye un antecedente directo de la transferencia multimodal, porque intenta construir un espacio común para ambas modalidades. No obstante, su objetivo es proporcionar una representación general para múltiples tareas y no estudiar de forma específica la clasificación contextual de palabras durante lectura natural.

Además de estos modelos generales, se ha estudiado el preentrenamiento orientado al habla. \textit{The Brain's Bitter Lesson} compara diferentes objetivos autosupervisados y concluye que las tareas genéricas pueden escalar mejor que los objetivos supervisados estrechamente vinculados a una etiqueta concreta \cite{jayalath2025bitter}. La mejora obtenida al aumentar la cantidad de datos sin etiquetar apoya la utilización de grandes corpus cerebrales antes del ajuste fino. Sin embargo, las ventanas utilizadas por este tipo de modelos suelen abarcar segundos y no minutos, por lo que no aprenden necesariamente a seleccionar dependencias lingüísticas de largo alcance.

En general, los modelos preentrenados de EEG han mejorado la reutilización entre sujetos y conjuntos, pero sus ventajas no son uniformes. El cambio de número de canales, frecuencia de muestreo, referencia y montaje puede exigir adaptadores específicos. Además, un modelo que funciona bien en sueño, epilepsia o imaginación motora no tiene por qué conservar la información necesaria para distinguir palabras. La evaluación sobre una tarea lingüística y frente a una inicialización aleatoria resulta, por tanto, imprescindible para comprobar que la transferencia aporta una mejora real.

\subsection{D'Ascoli et al. (MEG)}

El trabajo de d'Ascoli et al. no es un modelo preentrenado en sentido estricto. Se incluye en esta sección porque define el principal modelo supervisado de referencia para MEG-XL y formaliza la tarea de clasificación contextual de palabras que se adapta en este trabajo \cite{dascoli2024words}. La propuesta se entrenó sobre varios conjuntos EEG y MEG que, en conjunto, incluían cientos de participantes, diferentes idiomas y tareas de lectura y escucha.

Cada ejemplo se construye a partir de una ventana de tres segundos alineada con el inicio de una palabra. En primer lugar, un módulo de atención espacial combina los sensores teniendo en cuenta sus posiciones. Después se aplica una transformación dependiente del sujeto y una sucesión de bloques convolucionales dilatados. Este codificador produce una representación local para cada ventana. Las representaciones de todas las palabras de una oración se ordenan y se introducen en un Transformer bidireccional, de forma que la predicción de una posición pueda utilizar el contexto neuronal del resto de la frase.

El objetivo no consiste en predecir directamente una clase, sino en aproximar un \textit{embedding} de la palabra obtenido con T5. Para entrenar el espacio común se utiliza D-SigLIP, una modificación de SigLIP que evita penalizar como negativos dos ejemplos asociados a la misma palabra. Durante la evaluación, la representación predicha se compara por similitud coseno con un conjunto de recuperación. El resultado principal se expresa mediante precisión Top-10 equilibrada sobre las 250 palabras más frecuentes de cada conjunto.

La incorporación del Transformer contextual produce una mejora media cercana al 50\% respecto al codificador que procesa las palabras de forma aislada. El estudio también muestra que el rendimiento aumenta con la cantidad de datos de entrenamiento y con el promedio de varias repeticiones. Al comparar protocolos, MEG supera claramente a EEG y la lectura supera a la escucha cuando se presentan las mismas frases. Asimismo, se observa que dedicar más tiempo a cada participante puede resultar más útil que repartir el presupuesto de grabación entre muchos sujetos con pocas sesiones.

Pese a estos avances, el modelo necesita anotaciones lingüísticas y se entrena de forma supervisada desde cero para la tarea. Su contexto queda limitado a la oración presentada y no aprende previamente las regularidades de grabaciones continuas sin etiquetas. La capa dependiente del sujeto facilita el entrenamiento conjunto, pero complica la adaptación a una persona completamente nueva. Por último, la evaluación principal se realiza sobre un vocabulario frecuente y cerrado. Estos límites explican por qué el trabajo constituye una referencia sólida para el ajuste fino, pero no resuelve el problema de aprender con pocos datos del sujeto objetivo.

\subsection{MEG-XL}

Gran parte de este trabajo se sustenta en la existencia y éxito que ha tenido MEG-XL. Éste parte de la hipótesis de que el contexto útil de una señal cerebral no se limita a la duración de la unidad que se desea clasificar. La respuesta a una palabra puede depender de la frase, del discurso, del estado del participante y de propiedades estables de la sesión. Sin embargo, la mayoría de los modelos preentrenados procesa ventanas de entre medio segundo y treinta segundos. MEG-XL amplía esta escala y utiliza muestras continuas de dos minutos y medio, entre cinco y trescientas veces más largas que las empleadas por modelos anteriores \cite{jayalath2026megxl}.

El preentrenamiento se realiza sobre aproximadamente 300 horas de MEG y más de 800 participantes, combinando grabaciones de reposo, tareas sensoriomotoras y experimentos de lenguaje. La señal se remuestrea a 50 Hz y se tokeniza canal a canal mediante BioCodec. Este tokenizer reduce la dimensión temporal y produce códigos discretos mediante cuantización vectorial residual. A cada token se añaden representaciones de la posición del sensor y de su tipo, diferenciando magnetómetros y gradiómetros.

Sobre estos tokens se aplica un Transformer con atención \textit{criss-cross}. La atención temporal relaciona instantes dentro de un mismo sensor, mientras que la atención espacial combina sensores en un mismo momento. De esta forma se evita construir una matriz de atención completa sobre todos los canales y tiempos. Durante el entrenamiento se ocultan bloques temporales contiguos para todos los sensores y el modelo debe recuperar sus tokens. El enmascaramiento de bloques largos reduce la posibilidad de resolver la tarea mediante una interpolación local y obliga a utilizar información más distante.

Para la evaluación lingüística, MEG-XL se ajusta sobre LibriBrain, Armeni y MEG-MASC. La tarea sigue el planteamiento de d'Ascoli et al.: se procesan secuencias de ventanas alineadas con palabras y se predicen \textit{embeddings} T5 mediante D-SigLIP. Los experimentos muestran que el modelo preentrenado ofrece su mayor ventaja cuando hay pocos datos del sujeto objetivo. En algunos escenarios alcanza un rendimiento comparable al modelo supervisado utilizando alrededor de una hora en lugar de decenas de horas. Cuando se dispone de un registro muy profundo de una sola persona, el modelo supervisado puede recuperar la diferencia, lo que indica que el preentrenamiento sustituye principalmente a la falta de datos individuales.

El análisis de la longitud de contexto muestra además que las representaciones mejoran cuando el modelo ha sido preentrenado con ventanas más largas. No basta con proporcionar una secuencia extensa durante la inferencia: el modelo debe haber aprendido previamente a seleccionar qué regiones lejanas son relevantes. Las capas iniciales tienden a utilizar relaciones locales y las posteriores amplían progresivamente su alcance, lo que sugiere una integración jerárquica del contexto.

MEG-XL supone el antecedente más próximo al presente trabajo, pero mantiene varias limitaciones. El preentrenamiento y la evaluación se realizan con MEG, la tarea final se restringe a habla percibida y parte de los experimentos utiliza un vocabulario de 50 palabras. El modelo tampoco estudia la adaptación a montajes EEG, donde cambian la referencia, la geometría, el número de canales, la amplitud y la distribución del ruido. Aunque BioCodec fue entrenado originalmente con EEG, el resto de las representaciones contextuales se aprende sobre MEG. Por tanto, queda abierta la cuestión de si ese conocimiento temporal y espacial puede transferirse a EEG durante lectura natural.

\section{Limitaciones de los trabajos existentes}

La primera limitación común es la cantidad y distribución de los datos. Muchos conjuntos reúnen un número elevado de participantes, pero cada uno aporta menos de dos horas. Este formato permite estudiar generalización entre sujetos, aunque dificulta entrenar modelos profundos sobre patrones individuales. Los conjuntos de larga duración muestran que el rendimiento continúa creciendo al añadir horas por persona \cite{sato2024scaling,ozdogan2025libribrain,dascoli2024words}. Sin embargo, registrar decenas de horas para cada nuevo usuario no es una solución práctica, especialmente en una aplicación clínica.

La segunda limitación es la variabilidad entre dominios. Dos conjuntos EEG pueden diferir en referencia, número de electrodos, frecuencia de muestreo, impedancias, filtros y sistema de coordenadas. Al cambiar de EEG a MEG se modifica además la magnitud física registrada. Los módulos de atención espacial, las máscaras de sensores y los \textit{embeddings} de posición reducen parte de esta heterogeneidad, pero no garantizan que una representación aprendida en una modalidad conserve su utilidad en la otra. BrainOmni constituye un primer intento de unificación, aunque todavía no se ha demostrado una transferencia específica de un modelo MEG de contexto largo hacia clasificación de palabras con EEG.

La tercera limitación procede de las diferencias entre tareas. La actividad registrada durante habla producida contiene componentes motores y musculares que no aparecen durante escucha o lectura silenciosa. La lectura mediante presentación rápida de palabras facilita la alineación, pero no reproduce el comportamiento natural de los ojos. La lectura libre es más realista, aunque requiere seguimiento ocular y genera fijaciones de duración variable, regresiones y palabras omitidas. En consecuencia, un modelo entrenado con escucha continua no puede aplicarse directamente a lectura natural sin adaptar la construcción de ventanas y la interpretación de los eventos.

También existe una limitación temporal. Los modelos de clasificación tradicionales procesan intervalos aislados y descartan la información anterior y posterior. Los Transformers contextuales incorporan frases completas, pero siguen trabajando con secuencias preparadas para una tarea supervisada. Los modelos fundacionales, por su parte, suelen preentrenarse con ventanas cortas. MEG-XL demuestra que aprender con minutos de señal puede mejorar la transferencia, aunque todavía no se conoce qué parte de esta mejora corresponde a contexto lingüístico, estabilidad del sujeto, características de la sesión o regularidades del dispositivo.

La evaluación constituye otro punto crítico. Las métricas de generación pueden aumentar gracias a las probabilidades de un modelo de lenguaje, aunque la entrada cerebral sea poco informativa. El uso de \textit{teacher forcing}, la repetición de frases entre particiones o la selección de hiperparámetros sobre el conjunto de prueba pueden producir resultados excesivamente optimistas \cite{jo2024eegtext}. Por ello, se requieren líneas base con ruido, permutación de etiquetas, separación por estímulos y comparación entre la misma arquitectura inicializada aleatoriamente y preentrenada. En clasificación de palabras también debe indicarse el tamaño del vocabulario y utilizar métricas equilibradas para que las palabras frecuentes no dominen el resultado.

El vocabulario cerrado continúa siendo una restricción importante. La recuperación entre 50 o 250 palabras permite medir avances de forma controlada, pero está lejos de una comunicación abierta. Los mecanismos de reordenamiento e inserción mediante modelos de lenguaje amplían la cobertura \cite{jayalath2025unlocking}, aunque introducen el riesgo de que la coherencia lingüística oculte errores del decodificador cerebral. Antes de avanzar hacia una generación completamente abierta resulta necesario mejorar la representación neuronal y demostrar que cada incremento de rendimiento procede de la señal.

Por último, los trabajos actuales suelen estudiar por separado el preentrenamiento, la clasificación de palabras y la transferencia entre modalidades. Los modelos generales de EEG no se diseñan específicamente para lenguaje; el modelo contextual de d'Ascoli et al. necesita muchos datos etiquetados; y MEG-XL no ha sido evaluado con EEG. Esta separación define el hueco abordado en el presente trabajo. Lo que este trabajo propone es mantener la tarea y la arquitectura de clasificación contextual, incorporar conjuntos EEG de lectura y percepción, y comparar inicializaciones desde cero y preentrenadas. De este modo puede evaluarse si las representaciones de contexto largo aprendidas con MEG aportan una ventaja real al cambiar de modalidad o si, por el contrario, las diferencias físicas y experimentales producen transferencia negativa.
